\subsection{Cross-Architecture Measurement}
This study was exclusively conducted on an Intel Core i7-14700K CPU, which is a high-end desktop processor. Now, an open question is if the language hierarchy identifies here generalizes to other architectures. Processors like ARM, Ryzer, and Apple M-series etc. dominate the mobile and laptop market now, and differ from x86 architecture in terms of energy cost, branch prediction, and memory management. Future work can extend this experiment to a more diverse set of hardware platforms to see if the language hierarchy holds across different architectures. Additionally, the experiment can be repeated on a more recent generation of CPUs to see if the language hierarchy changes with newer hardware optimizations and energy efficiency improvements.


\subsection{Cloud Computing Energy Measurement}
All measurements in this study were collected on a local desktop machine where Intel RAPL is directly accessible via a kernel driver. Extending this methodology to public cloud environments like Amazon EC2, Google Cloud, or Microsoft Azure is both a challenge and an opportunity for future work. Cloud providers typically do not expose hardware-level energy measurements to users, and the virtualized nature of cloud in
